\capitulo{6}{Trabajos relacionados}

\section{Artículos científicos relacionados}

En este apartado se revisan los avances recientes en segmentación semántica aplicada a la detección de estructuras lineales, como lo son las sendas de fauna o los caminos rurales en imágenes aéreas o satelitales, especialmente en contextos ecológicos y de conservación. Se incluyen trabajos enfocados en el uso de aprendizaje profundo mediante redes neuronales convolucionales para identificar tales patrones lineales en terreno natural, con casos de estudio internacionales.

\begin{itemize}
\item\textbf{Remotely-Sensed Game Trails Are a Behavioral Footprint That Explains Patterns of Herbivore Habitat Use~\cite{Stears2025}}

Este estudio pionero analiza sendas de fauna en la sabana africana usando imágenes aéreas para inferir el uso de hábitat por herbívoros. Los autores cuantificaron la densidad de sendas mediante la fórmula: $\mathrm{m/ha}$ a partir de ortoimágenes de alta resolución y la modelaron frente a variables ambientales. Encontraron que la cobertura de vegetación leñosa es el mejor predictor: las mayores densidades de sendas ocurren con coberturas intermedias de matorral. Las sendas detectadas mostraron altas correlaciones con otros indicadores de uso de hábitat, como por ejemplo, los conteos de excrementos, validando su utilidad como huella del comportamiento de herbívoros a escala de paisaje. Este trabajo revela por primera vez que las sendas visibles en imágenes son un indicador integrador de los patrones de movimiento de fauna, proporcionando datos novedosos para ecología del paisaje y conservación.

\item\textbf{Human Inspired Deep Learning to Locate and Classify Terrestrial and Arboreal Animals in Thermal Drone Surveys~\cite{Backman2025}} 

Enfocado en monitoreo de fauna por drones, este artículo emplea técnicas de deep learning inspiradas en la percepción humana para detectar y clasificar animales terrestres y arborícolas en imágenes térmicas aéreas. Se desarrolló una red neuronal convolucional entrenada con vuelos de dron en Nueva Gales del Sur, Australia, para identificar especies crípticas como koalas y petauros en imágenes nocturnas. El modelo logró resultados robustos en la detección multiespecífica, demostrando el potencial de las \textit{CNN} para procesamiento autónomo de imágenes térmicas en programas de monitoreo de vida silvestre a gran escala.

\item\textbf{Automized Segmentation of Animal Paths in UAS Imagery Using Deep Learning — A Case Study from the Kruger National Park~\cite{Muller2025}} 

Este trabajo combina drones y aprendizaje profundo para mapear las rutas de paso de fauna en sabana africana. Se obtuvo mosaicos ortofotográficos RGB de alta resolución con dron en el Parque Kruger, en Sudáfrica y se entrenó una \textit{UNet} para segmentar automáticamente la red de sendas generadas por animales. Alcanzaron hasta un 83\% de precisión en la detección de sendas, demostrando la eficacia del modelo aún cuando muchas sendas son prácticamente invisibles al ojo humano. Un desafío destacado fue la generación de datos de entrenamiento de calidad: los autores implementaron un método semiautomático para etiquetar sendas, estimando su anchura a partir de datos de elevación fotogramétrica, lo que mejoró sustancialmente el rendimiento del modelo. Este estudio, fruto de la colaboración entre la Universidad de Würzburg y Parques Nacionales sudafricanos, muestra una herramienta escalable para investigar corredores de movimiento animal sin recurrir a collares GPS, complementando métodos tradicionales con mapas continuos de sendas.

\item\textbf{Mapping Trails and Tracks in the Boreal Forest Using LiDAR and Convolutional Neural Networks~\cite{McDermid2025}} 

En ecosistemas boreales de Canadá, este estudio mapea sendas y huellas en bosques boreales, tanto de fauna silvestre como senderos de vehículos creados por todoterrenos, empleando una combinación innovadora de \textit{LiDAR} de alta densidad y \textit{CNN}. Se utilizó un modelo \textit{U-Net} entrenado sobre modelos digitales de terreno (\textit{DTM}) de 10 cm de resolución obtenidos con \textit{LiDAR} montado en dron, complementado con datos \textit{LiDAR} aéreos de 50 cm. El modelo logró identificar automáticamente una extensa red de 2.829 km de senderos en un área de 59 $\mathrm{km}^2$ de bosque boreal, con alta exactitud, obteniendo: $F_1 \approx 0.74$--$0.77$. Se observó que las sendas se concentran especialmente en turberas y que a menudo aprovechan perturbaciones lineales humanas como parte de la red de movilidad. Este trabajo subraya la escasez previa de literatura sobre detección remota de sendas, demostrando el potencial de fusionar sensores activos \textit{LiDAR} con \textit{Deep Learning} para cartografiar rasgos lineales sutiles en ambientes naturales a gran escala, lo que puede informar tanto la gestión de fauna como la restauración de disturbios.

\item\textbf{Estimating Forest Biomass Using Small Footprint LiDAR Data: An Individual Tree-Based Approach That Incorporates Training Data~\cite{bortolot2005}} 

Aunque enfocado en biomasa forestal y no en sendas, se incluye este trabajo clásico por su relevancia metodológica en el uso de datos \textit{LiDAR} de alta resolución. Bortolot y Wynne desarrollaron un algoritmo a nivel de árbol individual para estimar biomasa aérea mediante \textit{small-footprint LiDAR}, demostrando la factibilidad de predecir variables ecológicas con alta precisión espacial. Aplicaron regresiones entre la densidad de retornos \textit{LiDAR} por árbol y mediciones de biomasa, logrando resultados prometedores. Si bien este estudio no trata sendas, sentó bases en el procesamiento de datos \textit{LiDAR} para estructuras ecológicas, técnica que, como muestran McDermid et al.~\cite{McDermid2025}, hoy facilita la detección de rasgos lineales sutiles, como sendas o cauces, en entornos naturales. Es un ejemplo temprano de cómo combinar teledetección detallada con algoritmos especializados puede extraer información ecológica no evidente mediante métodos tradicionales.
\end{itemize}

\section{Herramientas Visuales para Imágenes Satelitales y Aéreas}

En paralelo a los avances científicos, existen diversas herramientas de software que facilitan el acceso, visualización y análisis de imágenes aéreas y satelitales. A continuación se describen las más relevantes para proyectos de ecología del paisaje, excluyendo Google Earth Engine, ya que ha sido empleada en este TFG. Cada herramienta se caracteriza por su propósito, funciones clave, tipos de datos soportados y posibles usos en la detección de estructuras lineales naturales:

\begin{itemize}
\item\textbf{Interfaz de QGIS~\cite{qgis_web}} 

Un Sistema de Información Geográfica open source\footnote{\url{https://qgis.org/}} de escritorio, ampliamente utilizado por su potencia y flexibilidad. QGIS permite visualizar y procesar prácticamente cualquier formato ráster o vectorial, incluyendo extensiones como: GeoTIFF, JPEG2000, Shapefiles, etc..., donde se encuentran ortofotos aéreas y mosaicos satelitales. Ofrece herramientas de GIS avanzadas, siendo algunas de estas: la georreferenciación, el análisis de terreno y los filtros espaciales, y puede integrarse con datos remotos mediante complementos. Por ejemplo, el plugin \textit{QuickMapServices}~\cite{qgis_quickmapservices} brinda acceso directo a mapas base satelitales de diversas fuentes: Google, Bing, OSM, etc..., y el plugin \textit{Sentinel Hub}~\cite{sentinelhub_qgis_plugin} conecta con la plataforma Copernicus. 

QGIS es ideal para proyectos de conservación pues permite superponer capas y realizar análisis espaciales en un entorno unificado. Su capacidad de \textit{scripting} mediante Python y sus múltiples plugins facilitan la automatización de tareas como detección de caminos rurales o generación de mapas de densidad de sendas. Al ser software libre, es altamente personalizable y cuenta con una gran comunidad, con casos de uso en administraciones ambientales, universidades de España para gestión de hábitats y planificación territorial.

\item\textbf{USGS EarthExplorer~\cite{usgs_earthexplorer_web}}

Es el portal web de la agencia geológica estadounidense (USGS) para buscar, visualizar y descargar imágenes de satélite y fotografía aérea de múltiples fuentes. EarthExplorer funciona como una interfaz de descubrimiento de datos: el usuario delimita un área de interés en un mapa interactivo, define rangos de fechas y tipos de producto, como, por ejemplo, colecciones Landsat, Sentinel-2 y fotografías aéreas históricas, y el sistema devuelve las escenas disponibles. Incluye opciones avanzadas de filtrado, siendo algunas: por cobertura de nubes, ángulo solar, etc..., y permite examinar imágenes de baja resolución antes de ordenar la descarga. Es necesario un registro gratuito para obtener los datos. 

EarthExplorer da acceso a millones de imágenes de archivo desde los años 1960 en adelante, incluyendo datos de libre acceso como Landsat de niveles L1/L2, imágenes desclasificadas de satélites espía (CORONA), fotografías aéreas USDA, modelos digitales de elevación, entre otros. Su utilidad en proyectos ecológicos radica en compilar series temporales de imágenes históricas de un área. Un ejemplo aplicado a nuestro caso, podría ser el observar la apertura de nuevos caminos o sendas a lo largo de décadas y visualizar los cambios en cobertura vegetal asociados a herbívoros. También facilita obtener imágenes multiespectrales calibradas con las que aplicar algoritmos de segmentación entrenados en diferentes fuentes. Si bien EarthExplorer no permite un análisis geoespacial complejo por sí mismo, es un nodo central de acceso a datos de teledetección clásicos indispensable en flujos de trabajo reproducibles.

\item\textbf{Sentinel Hub EO Browser~\cite{sentinelhub_eobrowser_web}} 

Plataforma web desarrollada por Sinergise que ofrece una visualización instantánea de datos satelitales en la nube, con énfasis en las misiones Sentinel, del programa Copernicus, y otros datasets abiertos. EO Browser permite explorar y comparar imágenes en alta resolución de múltiples fuentes, como: Sentinel-1, -2, -3, Landsat 5-9, MODIS, entre otros, de forma interactiva. El usuario puede elegir un área simplemente navegando por el mapa, seleccionar un rango de fechas y filtrar por un porcentaje de nubes. La aplicación carga en segundos las imágenes disponibles cumpliendo esos criterios y las muestra en pantalla. Ofrece también visualizaciones personalizables, cmo pudieran ser las combinaciones de bandas espectrales, los índices de vegetación NDVI/NDWI predefinidos, las capas de falso color, etc..., y herramientas integradas para descargar imágenes en alta resolución, generar animaciones timelapse multitemporales, o incluso ejecutar scripts de procesamiento personalizados sobre los píxeles. 

EO Browser es una demo gratuita de las capacidades del servicio Sentinel Hub, por lo que no requiere al usuario instalar software ni manejar grandes volúmenes de datos localmente: todo el procesamiento se hace en servidores en tiempo real. Para proyectos de detección de sendas o caminos, EO Browser resulta muy útil en la fase exploratoria, ya que permite revisar rápidamente imágenes Sentinel-2 recientes de nuestro sitio de estudio, comparar épocas, identificar visualmente posibles patrones lineales y acotar períodos de interés. Si se necesita análisis más profundo, las imágenes encontradas se pueden descargar para tratamiento en QGIS u otro entorno. En resumen, EO Browser acerca el \textit{big data} satelital a cualquier usuario: con unos clics, se accede a décadas de imágenes globales, lo que empodera proyectos de conservación con información actualizada casi en tiempo real, ya que Sentinel-2 tiene una revisita de 5 días y los datos suelen estar disponibles horas después de su adquisición.

\item\textbf{Copernicus Browser~\cite{copernicus_browser_web}} 

Visor en línea oficial del programa Copernicus de la Agencia Espacial Europea, útil para acceder a los datos Sentinel de forma sencilla. Lanzado en 2023 dentro del Copernicus Data Space Ecosystem, este navegador gratuito y abierto presenta una experiencia muy similar a EO Browser, con enfoque en la facilidad de uso para el público general. Al ingresar, se muestra inmediatamente un mosaico sin nubes de Europa; el usuario puede navegar, hacer zoom y visualizar imágenes recientes prácticamente de cualquier lugar del mundo. Con un par de clics mediante el botón: “Visualize”, el sistema obtiene la última imagen Sentinel-2 sin nubes de la zona visualizada. Copernicus Browser integra también datos de Sentinel-1 (radar), Sentinel-3 (medio ambiente) y Sentinel-5P (atmósfera), permitiendo alternar entre diferentes misiones. Incorpora herramientas interactivas como vista 3D del terreno, medición de distancias/áreas, comparación temporal de imágenes, con: \textit{timelapse} y el deslizador antes-después y colecciones temáticas predefinidas, destacando fenómenos de interés, siendo algunos: la geología, la hidrología, los cambios urbanos, etc...

Un punto fuerte es que no requiere experiencia GIS: está diseñado para que incluso estudiantes o gestores no especializados exploren imágenes satelitales fácilmente. Esto resulta valioso en proyectos de biodiversidad donde se quiera involucrar a comunidades locales o personal de campo: por ejemplo, los guardabosques pueden abrir el navegador y ver rápidamente dónde han surgido nuevas sendas de ganado en una reserva, comparando con meses previos sin tener que procesar datos. Aunque no ofrece tantas opciones avanzadas como Sentinel Hub, ya que no se pueden ejecutar scripts personalizados, Copernicus Browser garantiza acceso abierto y streaming fluido de datos Copernicus, siendo excelente para visualización y educación ambiental. Para análisis detallados, los usuarios pueden registrarse y descargar los productos originales Sentinel desde la misma interfaz. En resumen, es una herramienta intuitiva y poderosa para poner imágenes satelitales al alcance de cualquier interesado en monitorear el territorio.

\begin{figure}
	\centering
	\includegraphics[width=0.7\linewidth]{Copernicus_Browser}
	\caption[Vista de la interfaz web de Copernicus Browser, mostrando un mosaico satelital en color real sobre Europa.]{}
	\label{fig:copernicus-browser}
\end{figure}


\item\textbf{Otras herramientas}

Además de las anteriores, cabe mencionar brevemente otras plataformas útiles según la finalidad: ESA SNAP, que consiste en un software de escritorio para procesado avanzado de imágenes de radar y ópticas, muy usado en investigación, Google Earth Pro, que permite la visualización 3D de imágenes históricas, aunque con limitaciones de resolución temporal en áreas naturales, NASA Worldview, un visor web de imágenes globales diarias, útil para ver rápidamente incendios, inundaciones u otros eventos, o visores regionales como el Centro de Descargas del CNIG en España. 

No obstante, para este proyecto centrado en sendas ecológicas, las herramientas ya descritas, QGIS, EarthExplorer, Sentinel Hub EO Browser y Copernicus Browser, cubren un espectro completo: desde el análisis GIS exhaustivo offline, pasando por la búsqueda masiva de datos brutos, hasta la visualización interactiva online. Integrar estas herramientas en el flujo de trabajo permitirá identificar, segmentar y validar las estructuras lineales de interés con mayor eficacia y rigor.
\end{itemize}
